\setcounter{section}{-1}
\section{Allgemeine Voraussetzungen und Konventionen}

Es werden umfassende Kenntnisse der 
\emph{linearen Algebra},
\emph{Funktionentheorie} und grundlegende Begriffe der \emph{Analysis auf Mannigfaltigkeiten} -- insbesondere \emph{topologische Grundbegriffe} -- vorausgesetzt. Außerdem sollte der Begriff einer \emph{Gruppenoperation} bekannt sein. 
\begin{notation}[Indexmenge]
    Falls nichts anderes angegeben, bezeichnet \(I\) eine beliebige Indexmenge.
\end{notation}
\begin{notation}[Element+Menge]
	Sei \(X\) eine Menge, \(A,B\subseteq X\) Teilmengen und \(\cdot\) eine Verknüpfung von Elementen in \(X\). Dann schreiben wir
	\[x\cdot A := \{x\cdot a\mid a\in A\}, \quad A\cdot B := \{a\cdot b \mid a\in A, b\in B\}.\]
\end{notation}

\pagebreak
